%-----------------------------------------------------------------------------------
%Bu latex dökümanı AUBB için tez yazım dökümanı olarak
%hazırlanmıştır. Sürekli olarak değişiklikler yapılmaktadır.
%En son yazılmış olan dosyaları kullanmaya özen gösteriniz.
%Son değişiklik tarihi: 05 Temmuz 2022 dür.
%Son değişiklik yapan: Mustafa K. Yıldız.
%%%
\documentclass[12pt, twoside, a4paper]{aubb}

\makeatother

%%%%%%%%%%%%%%%%%%%%%%%%%%%%%%%%%%%%%%%%%%%%%%%%%%%%%
%%%%%%%%%%%%%%%%%%%%%%%%%%%%%%%%%%%%%%%%%%%%%%%%%%%%%
%%% Tez bölümü bu kısımdan başlar
\begin{document}

\newgeometry{hmarginratio=1:1}    %% make layout symmetric
\include{Bölümler/0-Kapak.tex}
\restoregeometry 

\pagenumbering{roman} %% Bu komut sayfa numarasının tipini belirler.
\setcounter{page}{3}
%%%%%%%%%%%%%%%%%%%%%%%%%%%%%
%%%%%%%%%%%%%%%%%%%%%%%%%%%%%%%%%%%%%%%%%%%%%%%%%%%%%%

\include{Bölümler/01-Juri.tex}
\cleardoublepage


%%  Özet kısmı (Türkçe).

\include{Bölümler/1-Özet.tex}
\cleardoublepage

%%% Özet kısmı (İngilizce). 

\include{Bölümler/2-Özet-ingilizce.tex}
\cleardoublepage

%%% Teşekkür kısmı
\include{Bölümler/3-Teşekkür.tex}
\cleardoublepage

%%%%%%%%%%%%%%%%%%%%%%%%%%%%%%%%%%%%%%%%%%%%%%%%%%%%%
%%% Şekillere ve Tablolara verilen numaralar sayesinde sadece aşağıdaki komutlarla şekil ve tablo dizini oluşturulabilir.
%%%%%%%%%%%%%%%%%%%%%%%%%%%%%%%%%%%%%%%%%%%%%%%%%%%%%


%%% İÇİNDEKİLER
\tableofcontents
\cleardoublepage

%%% ŞEKİLLER
\addcontentsline{toc}{section}{ŞEKİLLER DİZİNİ} 
\listoffigures
\cleardoublepage

%%% TABLOLAR
\addcontentsline{toc}{section}{ÇİZELGELER DİZİNİ}
\listoftables
\cleardoublepage


%%%%%%%%%%%%%%%%%%%%%%%%%%%%%%%%%%%%%%%%%%%%%%%%%%%%%
%%%% Bundan sonra sayfa sayıları 1,2,3 arap gösterimi ile olacak.
\pagenumbering{arabic}
%%%%%%%%%%%%%%%%%%%%%%%%%%%%%%%%%%%%%%%%%%%%%%%%%%%%%
%%% Artık tezinizin bölümlerini çağırmaya başlayabilirsiniz.
%%%%%%%%%%%%%%%%%%%%%%%%%%%%%%%%%%%%%%%%%%%%%%%%%%%%%

\include{Bölümler/Bölüm1.tex}

\cleardoublepage

\include{Bölümler/Bölüm2.tex}

\cleardoublepage

\include{Bölümler/Bölüm3.tex}

\cleardoublepage

\include{Bölümler/Bölüm4.tex}

\cleardoublepage

\include{Bölümler/Bölüm5.tex}

\cleardoublepage

\include{Bölümler/Bölüm6.tex}

\cleardoublepage

%%%%% Tez Sonu: Özgeçmiş

\include{Bölümler/4-Özgeçmiş.tex}

\end{document}
